% =============================================================================
% Section 15 — Future Work
% =============================================================================
\section{Future Work}
\label{sec:future}

This section outlines the planned evolution of \pebble{} beyond v0.1, grouped
by impact and estimated implementation complexity.

\subsection{Hybrid Logical Clocks (HLC)}

\textbf{Motivation:} Replace physical wall-clock timestamps with Hybrid Logical
Clocks~\cite{kulkarni2014hlc} to capture causality between events, even under
clock skew.

\textbf{Approach:} Each node maintains an HLC state $(l, c)$ where $l$ is the
maximum of the local physical clock and the highest $l$ seen in received
messages, and $c$ is a counter that breaks ties at the same $l$ value. The HLC
advances monotonically and satisfies $l \geq \mathrm{now}()$ at all times,
ensuring that causally later events always receive larger HLC values.

\textbf{Impact:} The LWW comparison function (\cref{def:lww-order}) would
replace $m.\mts$ with the HLC value, eliminating clock-skew-induced inversions.
The mutation wire format would add an HLC component alongside the physical
timestamp. Backward compatibility requires a protocol version bump.

\subsection{Sparse Version Vectors}

\textbf{Motivation:} Remove the contiguous-sequence assumption to support
out-of-order mutation delivery and multi-path replication.

\textbf{Approach:} Replace the high-water mark integer $V[n]$ with a compact
set representation of received sequence numbers (e.g., a sorted interval list
or a roaring bitmap~\cite{riak2010}). The \texttt{getMissingRanges} function
would compute gaps in the received set rather than the single suffix above the
watermark.

\textbf{Impact:} Memory overhead per node increases by a constant factor. Gossip
requests become slightly larger (a list of gaps rather than a single range).
Enables future UDP-based or multi-path gossip without correctness risk.

\subsection{History Compaction and Tombstone GC}

\textbf{Motivation:} Bound memory and on-disk storage by discarding history
entries for mutations that are no longer needed for gossip sync.

\textbf{Approach:} Periodically take a snapshot and promote it to a
\emph{stable checkpoint}. All history entries for mutations covered by the
checkpoint can be discarded. Gossip peers that have not yet received those
mutations must be directed to fetch the snapshot instead of individual mutations
(a form of state transfer).

Tombstone GC requires additional coordination: a tombstone can only be discarded
if all nodes in the cluster have seen it. A simple mechanism is a cluster-wide
minimum version vector: $V_{\min}[n] = \min_{i} V_i[n]$ over all nodes. History
entries with sequence $\leq V_{\min}[n]$ are safe to discard.

\textbf{Impact:} Bounds $|H|$ to $\Oof{K}$ after compaction (one entry per key).
Requires nodes that fall far behind to perform a full state transfer rather than
incremental gossip.

\subsection{Dynamic Cluster Membership}

\textbf{Motivation:} Allow nodes to join or leave the cluster without restarting
all nodes or manually editing \texttt{cluster.json}.

\textbf{Approach:} Represent the cluster membership as a CRDT (e.g., a 2P-Set
or OR-Set) replicated through the gossip protocol. A joining node broadcasts a
signed join request; existing nodes add it to the membership CRDT and begin
gossiping with it. A leaving node broadcasts a signed leave message. The
moderator signs membership changes to prevent rogue nodes from injecting
themselves.

\textbf{Impact:} Eliminates the need for out-of-band cluster reconfiguration.
Requires all nodes to run a compatible membership CRDT implementation.

\subsection{Multi-Key Transactions}

\textbf{Motivation:} Provide atomic read-modify-write and multi-key update
semantics for applications that require them.

\textbf{Approach:} Two candidate approaches:
\begin{enumerate}[leftmargin=1.5em]
  \item \textbf{Optimistic concurrency control (OCC):} The client reads a set
        of keys with their current version vectors, applies changes, and submits
        a conditional write that is rejected if any observed version has changed.
        Conflicts are retried by the client.
  \item \textbf{Saga pattern:} Multi-key updates are composed as a sequence of
        compensating single-key mutations. Partial failures leave a detectable
        intermediate state that can be rolled back.
\end{enumerate}
Full distributed serialisable transactions would require a consensus layer
(contradicting the leaderless design goal) and are not planned.

\subsection{TLS Transport Encryption}

\textbf{Motivation:} Protect mutation contents and client credentials from
on-path eavesdroppers.

\textbf{Approach:} Wrap both the gossip port and the client port in TLS 1.3.
Peer authentication during the TLS handshake can reuse the existing Ed25519
keys, eliminating the need for a separate PKI (mTLS with self-signed certs
derived from the Ed25519 identity). The application-level handshake
(\cref{sec:nodeauth}) could be simplified or merged with the TLS handshake.

\subsection{Configurable Gossip Fanout}

\textbf{Motivation:} Reduce convergence time for large clusters by initiating
$f > 1$ gossip sessions per interval.

\textbf{Approach:} Expose a \texttt{gossipFanout} parameter in
\texttt{cluster.json}. The \texttt{GossipEngine} selects $f$ random healthy
peers per interval and runs sessions concurrently. Expected convergence time
drops to $\Theta(\log_f N)$ rounds, as each round covers an $f$-fold more ground.

\subsection{Read Repair and Anti-Entropy}

\textbf{Motivation:} Detect and repair state divergence that persists beyond
the normal gossip window (e.g., due to prolonged partitions or log corruption
on a specific node).

\textbf{Approach:} A periodic anti-entropy pass computes a Merkle tree over
the current state and compares digests with all peers. Divergent subtrees are
repaired by re-exchanging the affected mutations. This provides a background
consistency guarantee independent of the gossip fanout.

\subsection{Read Consistency Modes}

\textbf{Motivation:} Offer stronger read guarantees for applications that need
them, without abandoning the leaderless design for writes.

\textbf{Approach:} Two optional read modes:
\begin{description}[leftmargin=1.5em]
  \item[Read-your-writes:] After a write at node $n$, route subsequent reads
        for the same key to $n$ until the version has propagated.
  \item[Monotonic reads:] Track a client-side version vector; reject reads from
        nodes whose vector does not dominate the client's last-observed vector.
\end{description}
Both modes are implementable without consensus and add only client-side
bookkeeping overhead.
